%% confusable_dd_optimization.tex
%%
%% Mathematical specification for Approach 2: finding optimal encoding directions
%% by maximizing the Lasserre SOS/SDP lower bound on the minimum confusable pair size.
%%
%% Companion to confusable_dd_search.py (proof-of-concept implementation for k=4).
%% June 2026.

\documentclass[12pt,a4paper]{article}

\usepackage{amsmath,amssymb,amsthm}
\usepackage{hyperref}
\usepackage{geometry}
\usepackage{listings}
\usepackage{xcolor}
\geometry{margin=2.5cm}

\newtheorem{definition}{Definition}
\newtheorem{theorem}{Theorem}
\newtheorem{proposition}{Proposition}
\newtheorem{remark}{Remark}
\newtheorem{assumption}{Assumption}

\lstset{
  language=Python,
  basicstyle=\small\ttfamily,
  keywordstyle=\bfseries,
  commentstyle=\color{gray},
  breaklines=true,
  frame=single,
  aboveskip=6pt,
  belowskip=6pt,
}

\title{Finding Optimal Encoding Directions via SOS/SDP Certificates}

\author{Lester (problem), Claude (approach + implementation)}
\date{June 2026}

\begin{document}
\maketitle

\begin{abstract}
This document formalizes Approach 2 for finding the sequence of optimal encoding
directions $DD(i)$ for $i = 1, 2, \ldots$ in dimension $k$ that maximise the
minimum confusable pair size $f(DD(i))$.  The method reformulates the problem as
an outer optimization over directions, with an inner oracle that certifies
lower bounds via the Lasserre moment hierarchy.  For $k=4$ and $s=2$, this is
concretely implementable using exact rational arithmetic and symbolic
verification.  The approach is rigorous (no tolerances, no measure-zero
arguments) but slow.
\end{abstract}

\tableofcontents
\newpage

%% -----------------------------------------------------------------------
\section{Problem statement}
\label{sec:problem}

Given dimension $k$ and number of directions $i$, we seek to find a set of
$i$ encoding directions
$$D^* = (d_1^*, \ldots, d_i^*) \in (S^{k-1})^i$$
that maximises the minimum confusable multiset size:
$$f(D^*) = \max_{D \in (S^{k-1})^i} \min\{\, s : \text{no confusable pair of size } < s \text{ exists for } D \,\}.$$

This is an outer optimisation (find best $D$) with an inner oracle (determine
$f(D)$ for fixed $D$).  We cannot compute $f(D)$ exactly in general, but we can
compute a **lower bound** via the Lasserre certificate.

\begin{remark}
  We seek $DD(i)$ with $f(DD(i))$ rigorous (certified) — not approximate.
\end{remark}

%% -----------------------------------------------------------------------
\section{The inner oracle: Lasserre certificates for fixed directions}
\label{sec:inner-oracle}

\subsection*{Setup}

Fix encoding directions $D = (d_1, \ldots, d_m)$ with $m = i$.  The question
``does a confusable $s$-element pair exist?'' is equivalent to: does the
polynomial system
\begin{equation}
  \label{eq:confusable-system}
  \sum_{v \in X} (v \cdot d_j)^r = \sum_{v \in Y} (v \cdot d_j)^r,
  \quad j=1,\ldots,m,\; r=1,\ldots,s,\quad X \neq Y
\end{equation}
have a solution?

After centring both $X$ and $Y$ at the origin (so $\sum_{v \in X} v = \sum_{v \in Y} v = 0$),
this reduces to $(s-1)m$ equations in $2k(s-1)$ variables (the centered coordinates).

\subsection*{The SOS certificate (degree-4 case, $s=2$)}

For $s=2$, the system reduces to power sums of projections agreeing for $r=1, 2$.
The $r=1$ condition forces equal centroids (achieved by centring); the $r=2$ condition is
\[
  \sum_{v \in X} (v \cdot d_j)^2 = \sum_{v \in Y} (v \cdot d_j)^2 \quad \forall j.
\]

Define the polynomial
\[
  f_D(\vec{A}, \vec{B})
  = \sum_{j=1}^m (\vec{A} \cdot d_j)^2 (\vec{B} \cdot d_j)^2,
  \qquad
  g_D(\vec{A}, \vec{B}) = \|\vec{A}\|^2 \|\vec{B}\|^2,
\]
where $\vec{A}, \vec{B} \in \mathbb{R}^k$ represent the two elements of $X$ and $Y$ (centred).

A confusable pair of size 2 exists if and only if $f_D(\vec{A}, \vec{B}) = 0$ for some
$(\vec{A}, \vec{B}) \neq (0, 0)$, i.e. if $\min f_D / g_D = 0$.

\begin{definition}
  The \emph{s=2 SOS certificate} for directions $D$ is the maximum real number
  $\lambda^*(D)$ such that there exist a PSD matrix $Q$ and polynomials $p_\gamma$
  (indexed by monomials $\gamma$ of degree $\leq 4$ in $2k$ variables) satisfying
  \[
    f_D - \lambda^*(D) \cdot g_D = \sum_\gamma p_\gamma(\vec{A}, \vec{B})^2.
  \]
\end{definition}

\begin{proposition}
  \label{prop:sos-lower-bound}
  If $\lambda^*(D) > 0$, then $f_D(\vec{A}, \vec{B}) \geq \lambda^*(D) \cdot g_D(\vec{A}, \vec{B})$
  for all $(\vec{A}, \vec{B})$.  Hence, no non-trivial confusable pair of size 2 exists,
  and $f(D) \geq 3$.
\end{proposition}

If $\lambda^*(D) = 0$, the certificate is inconclusive (a confusable size-2 pair may or may not exist).

\subsection*{Computation via SDP}

The SOS certificate is computed by solving a semidefinite programme (Section
\ref{sec:sdp-formulation} below).  The key structure:
- **Decision variables**: the Gram matrix $Q$ and the scalar $\lambda$
- **Objective**: maximise $\lambda$
- **Constraints**: $Q \succeq 0$ (PSD) and polynomial-identity constraints (linear in $Q, \lambda$)

\subsection*{Generalisation to $s \geq 3$ via Lasserre hierarchy}

For $s \geq 3$, the SOS approach extends via the Lasserre moment hierarchy (relaxation
level $d$).  The moment matrix grows as $\binom{2k+d}{d}$, and higher $d$ gives tighter
(larger) lower bounds.  For this document we focus on $s=2$ (a single fixed SDP) and
note that $s \geq 3$ requires larger-scale SDPs.

%% -----------------------------------------------------------------------
\section{SDP formulation for k=4, s=2}
\label{sec:sdp-formulation}

For concreteness, we work in $k=4$ (dimension of vectors).  The directions $d_j$ are
4-dimensional unit vectors.

\subsection*{Variables and monomials}

The 2k=8 variables are the components of $(\vec{A}, \vec{B})$:
$$A = (A_1, A_2, A_3, A_4), \quad B = (B_1, B_2, B_3, B_4).$$

For degree $\leq 2$ monomials in 8 variables, there are $\binom{8+2}{2} = 45$ monomials.
The vector $v = (1, A_1, \ldots, A_4, B_1, \ldots, B_4, A_1^2, A_1 A_2, \ldots, B_4^2)$
(all monomials of degree $\leq 2$) has length 45.

For degree $\leq 4$ monomials, there are $\binom{8+4}{4} = 495$ monomials.

\subsection*{Gram matrix and polynomial identity}

We seek a $45 \times 45$ PSD matrix $Q$ (the Gram matrix) such that
\[
  v^T Q v = f_D(A, B) - \lambda \cdot g_D(A, B)
\]
as a polynomial identity.  Expanding both sides and matching coefficients of each
degree-4 monomial $\gamma$:
\[
  \sum_{\substack{\alpha + \beta = \gamma \\ |\alpha|, |\beta| \leq 2}}
  Q_{\text{idx}(\alpha), \text{idx}(\beta)}
  = f_{D,\gamma} - \lambda \cdot g_{D,\gamma},
\]
where $\text{idx}(\alpha)$ is the index of monomial $\alpha$ in the vector $v$, and
$f_{D,\gamma}, g_{D,\gamma}$ are the known coefficients of $f_D$ and $g_D$ at monomial $\gamma$.

\subsection*{The SDP}

\begin{align}
  \text{maximise} \quad & \lambda \\
  \text{subject to} \quad
    & Q \succeq 0, \\
    & v^T Q v = f_D - \lambda g_D \quad \text{(polynomial identity)}.
\end{align}

There are 495 linear constraints (one per degree-4 monomial) and 1 PSD constraint.

\subsection*{Solving with CVXPY + numerical verification}

We use CVXPY with a solver (SCS or MOSEK) to compute a numerical solution
$(Q_*, \lambda_*)$.  The solution is approximate (floating-point), but we verify it
rigorously (Section \ref{sec:verification}) before accepting it as a certificate.

%% -----------------------------------------------------------------------
\section{Rigorous verification without tolerances}
\label{sec:verification}

After the SDP solver returns a candidate $(\lambda_*, Q_*)$, we verify that it is a
valid SOS certificate using **symbolic computation and interval arithmetic**.

\subsection*{Verification procedure}

\begin{algorithm}
\textbf{Input:} directions $D$, candidate $(\lambda_*, Q_*)$ from SDP solver, dimension $k=4$.

\begin{enumerate}
  \item Construct $f_D, g_D$ as symbolic polynomials in SymPy.

  \item Construct the polynomial difference $P = f_D - \lambda_* \cdot g_D$ symbolically.

  \item Compute the Gram matrix product $v^T Q_* v$ symbolically (as a polynomial in
    $A_1, \ldots, B_4$).

  \item Check if $P = v^T Q_* v$ **exactly** (as polynomials in SymPy).  If yes, return
    \texttt{VERIFIED}.

  \item If not exactly equal (likely due to solver error), use interval arithmetic:
    expand both sides, compare coefficients with interval bounds, check if all intervals
    contain zero. If yes, the SOS certificate holds approximately and we can tighten
    $\lambda$ slightly (via a rounding procedure) to get an exact certificate.

  \item If verification fails, return \texttt{UNVERIFIED} and discard $(\lambda_*, Q_*)$.
\end{enumerate}
\end{algorithm}

\begin{remark}
  The key advantage: we never compare floating-point numbers directly. We either verify
  the identity exactly (symbolically) or we use **rigorous interval arithmetic** to
  certify the bounds. No $\epsilon$ tolerances on dot products or eigenvalue checks.
\end{remark}

%% -----------------------------------------------------------------------
\section{Parameterisation of directions}
\label{sec:parameterization}

To make the outer search finite, we parameterise unit vectors in $\mathbb{R}^k$ using
rational coordinates.

\subsection*{Rational points on the sphere}

\begin{definition}
  A \emph{rational unit vector} in $\mathbb{R}^k$ is a vector $d = (d_1, \ldots, d_k)$
  where each $d_i \in \mathbb{Q}$ and $\|d\|^2 = 1$ (exactly).
\end{definition}

For $k=4$, we can represent rational unit vectors as:
\[
  d = \frac{(q_1, q_2, q_3, q_4)}{\sqrt{q_1^2 + q_2^2 + q_3^2 + q_4^2}},
  \quad q_i \in \mathbb{Z}.
\]

The denominator $\sqrt{\sum q_i^2}$ is an algebraic number (the square root of an integer).
Representing and comparing such vectors exactly requires **algebraic arithmetic** (polynomial
rings over $\mathbb{Q}$, represented via SymPy's `sqrt` and `Rational`).

\subsection*{Practical strategy: bounded-numerator search}

To keep the search finite and tractable:

\begin{enumerate}
  \item Fix a maximum numerator $N$ (e.g., $N=20$).

  \item Enumerate all integer vectors $(q_1, \ldots, q_k)$ with $|q_i| \leq N$, $q_i \neq 0$.

  \item Normalize: $d = (q_1, \ldots, q_k) / \|q\|$ where $\|q\| = \sqrt{q_1^2 + \cdots + q_k^2}$.

  \item Store $d$ symbolically as a SymPy vector with algebraic coordinates.

  \item For each $i$-tuple of directions $(d_1, \ldots, d_i)$, compute the SDP and verify
    the certificate.

  \item Track the maximum $\lambda^*$ achieved.
\end{enumerate}

For $k=4$ and $N=20$: there are roughly $(2 \cdot 20)^4 / 6 \approx 100,000$ distinct unit
directions (accounting for duplicates after normalization). The search over all $i$-tuples
is combinatorially large ($\sim 100,000^i$), but we use heuristics to prune (e.g., only try
``interesting'' subsets: symmetric configurations, random samples, etc.).

\subsection*{Algebraic vs. floating-point coordinates}

When constructing $f_D$ and $g_D$ symbolically, the coefficients will involve square roots
(from the normalization).  SymPy can handle this, but arithmetic becomes slow.

**Alternative (faster):** Pre-normalize integer vectors $(q_1, \ldots, q_k)$ to float
coordinates $d = q / \|q\|_{\text{float}}$, compute the SDP with floats, then verify
the result.  This is not slower than floating-point but sacrifices the algebraic exactness.
It is still rigorous because the **verification** step (Section \ref{sec:verification})
is symbolic and uses interval arithmetic.

%% -----------------------------------------------------------------------
\section{Grid search and optimisation}
\label{sec:search}

\subsection*{Exhaustive search for small i}

For $i \leq 6$: enumerate all $i$-tuples $(d_1, \ldots, d_i)$ of unit directions from
the bounded-numerator set, compute $\lambda^*(D)$ for each, return $\arg\max \lambda^*$.

For $i=7, 8$: the search space is too large ($ \sim 10^{28}$); use sampling strategies below.

\subsection*{Sampling strategies for large i}

\begin{enumerate}
  \item **Symmetric configurations**: regular polytopes (e.g., 24-cell in $\mathbb{R}^4$
    has 24 vertices), spherical designs, equiangular tight frames.  Compute $\lambda^*$ for
    each and track best.

  \item **Random sampling**: draw random $i$-tuples of unit directions (uniformly on
    $(S^{k-1})^i$), compute $\lambda^*$, track best over many trials.

  \item **Local optimisation**: starting from a promising configuration, use gradient-based
    methods to improve it.  The gradient of $\lambda^*(D)$ w.r.t. $D$ can be computed from
    the SDP dual solution (envelope theorem).

  \item **Heuristic pruning**: discard $i$-tuples where some pair of directions is nearly
    parallel (likely to give small $\lambda^*$).
\end{enumerate}

\subsection*{Output and guarantees}

For a given $i$, the output is:
\begin{itemize}
  \item Directions $DD(i) = (d_1, \ldots, d_i)$
  \item Certified lower bound $\lambda^* > 0$ (if verified)
  \item Conclusion: $f(DD(i)) \geq 3$ (no size-2 confusable pair).
\end{itemize}

The true minimum confusable pair size $f(DD(i))$ may be $\geq 3$; we are only proving
a lower bound.  To determine if $f(DD(i)) = 3$ exactly, one would need to (a) search for
$s=3$ confusable pairs numerically, and (b) use higher Lasserre levels to tighten the bound.

%% -----------------------------------------------------------------------
\section{Implementation notes}
\label{sec:implementation}

The proof-of-concept code `confusable_dd_search.py` implements this approach for $k=4$.

\subsection*{Key functions}

\begin{enumerate}
  \item \texttt{sos\_sdp\_for\_directions(D, s=2)}: constructs and solves the SDP for
    given directions.  Returns $(\lambda_*, Q_*, \text{status})$.

  \item \texttt{verify\_sos\_certificate(D, lambda\_star, Q\_star)}: verifies the
    certificate using SymPy.  Returns \texttt{True/False}.

  \item \texttt{normalize\_to\_unit\_rational(q)}: normalizes an integer vector to a
    rational unit vector.  Returns symbolic SymPy vector.

  \item \texttt{enumerate\_bounded\_directions(k, N)}: generates all rational unit vectors
    with integer numerators $\leq N$.  Yields generators (lazy) to avoid memory overload.

  \item \texttt{search\_optimal\_directions(k, i, search\_type='exhaustive'|'random'|'symmetric')}:
    searches for the best $i$-tuple and returns $(DD, \lambda^*, \text{verified})$.
\end{enumerate}

\subsection*{Computational complexity}

- For each SDP: $O(495^3)$ (dense matrix operations on $495 \times 495$ moment matrix).
  With a sparse SDP solver, can be reduced.
- For each verification: $O(\text{poly}(k, s) \cdot \text{deg}(\text{SymPy algebra}))$.  Slow
  but exact.
- Grid size: $O((2N)^{k \cdot i})$ in worst case; reduced by symmetry and sampling.

For $k=4, i=4$: exhaustive search over all 4-tuples of directions with $|q_i| \leq 20$ is
$ \sim 10^{13}$ operations (slow but feasible with optimisation).  For $i \geq 7$: use
sampling.

\subsection*{Robustness and debugging}

- All SymPy expressions are pretty-printed for human inspection.
- Failed verifications are logged with the candidate $(\lambda_*, Q_*)$ for diagnosis.
- Interval-arithmetic bounds are tightened iteratively if a certificate is close to
  verification but fails numerically.

%% -----------------------------------------------------------------------
\section{Limitations and future work}
\label{sec:limitations}

\subsection*{Known limitations}

\begin{enumerate}
  \item **Limited to $s=2$** in the current proof-of-concept.  Extension to $s \geq 3$
    requires Lasserre level $d \geq 2$, which scales poorly in dimension.

  \item **Grid search is finite but potentially incomplete**.  We may miss the true
    optimal $DD(i)$ if it lies outside the bounded-numerator set.  However, we can increase
    the bound $N$ to search larger regions.

  \item **Verification is slow**.  SymPy symbolic computation on large polynomials is
    expensive.  For faster verification, use interval arithmetic alone (which is cheaper
    but less certain).

  \item **No negative results**.  If we fail to find a direction set with $\lambda^* > 0$,
    we cannot conclude that none exists — only that our search was incomplete.  To prove
    non-existence would require exhaustive search or a theoretical argument.
\end{enumerate}

\subsection*{Future extensions}

- Lasserre hierarchy for $s \geq 3$ (higher relaxation levels).
- Gradient-based outer loop to refine promising direction sets.
- Exact algebraic optimization (Gröbner bases) instead of numerical SDP solvers, for
  complete rigor without floating-point.
- Integration with the MinimalConfusableSets hypercube code to compare upper bounds.

%% -----------------------------------------------------------------------
\section{Relation to other approaches}
\label{sec:comparison}

\subsection*{vs. Approach 1 (exploit M < k structure)}

Approach 1 focuses on disproving L-matrix cancellations by hand for special cases.
Approach 2 is more systematic: it provides a machine-verifiable lower bound for any
direction set, without case analysis.

\subsection*{vs. direct polynomial system solving}

One could try to directly find confusable pairs for a given $D$ by solving the
polynomial system (\ref{eq:confusable-system}) with Gröbner bases or resultants.
This is more general but slower and gives only a yes/no answer, not a certified
lower bound.  Approach 2 complements this: if the polynomial system has no small
solution, the SOS certificate proves it rigorously.

\subsection*{vs. numerical search (confusable\_search.py)}

Numerical search (from the companion document `confusable_sdp.tex`) gives empirical
evidence that confusable pairs exist or don't.  Approach 2 provides *certified* evidence
via a mathematical proof (the SOS certificate).

\end{document}
